\section{Scheduling the first positive answer}
\label{sec:scheduling}

The maximum-depth algorithm must coordinate discovery with the heights
of its child codes. Expected depth poses a different scheduling need:
caps containing likely targets should be tested early, which leads to
weighted min-sum set cover (MSSC). We use the supplied caps as the
covering sets, with a singleton cap included whenever a supplied test
realizes it.

In the decision-tree application, a positive hull answer exits the
all-negative route and reveals that the target lies in a small cap.
The all-zero pattern follows that route through every test; once the
schedule has covered every positive pattern, the all-zero class is the
remaining outcome. An earlier singleton inference shortens its route.
Its expected-cost contribution is therefore bounded by its weight
times the schedule length. Accordingly, our ordering bound tracks both
the sum of positive-pattern first-cover times and the length paid on
the all-negative route.

\begin{symboltable}
$U,\ \pi_h,\ p_0$ & Coverable positive-pattern demands, their positive weights, and the weight of a target found only after the schedule finishes.\\
$t_Q(h),\ \MSSC(Q)$ & First-cover position of demand $h$ in ordering $Q$ and the sum $\sum_h\pi_h t_Q(h)$.\\
$k_0$ & Minimum number of candidate tests needed to cover all demands in $U$.\\
$\beta,\ \Phi_\beta,\ Z_\beta$ & Geometric budget factor and the first-cover and total-length accounting functions.\\
$r_Q(x),\ G_u$ & Comparison-order mass with first-cover time greater than $x$, and the constructed schedule for shift $u$.\\
\end{symboltable}

Weighted MSSC has a classical greedy four-approximation~\cite{flt2004}.
For circular intervals, we obtain a stronger guarantee by first solving
the maximum-coverage problem exactly for each cardinality budget.
Ordering these budget-optimal sets in geometric batches then controls
both first-cover cost and total length.

This batch method extends minimum cover by optimizing covered weight
at every integer budget. The interval structure makes each optimization
tractable: once the previous selected position is known, the next
position determines the newly covered demands.

\begin{lemma}[Exact budgeted interval coverage]\label{lem:budget-coverage}
For positively weighted interval or circular-interval demands on
a finite candidate set, a set of at most $k$ candidates hitting
maximum demand weight is computable in polynomial time.
\end{lemma}
\begin{proof}
Algorithm~\ref{alg:interval-cover} gives both interval oracles.
For ordinary intervals, number the candidates $1,\ldots,M$.
If the previous selected candidate is $p$ and the next is $q>p$,
the newly hit demands are the intervals containing $q$ whose left
endpoint is greater than $p$. An interval containing
$q$ and an earlier selected point must also contain $p$.
Conversely, any interval containing $q$ with left endpoint greater
than $p$ contains no earlier selected candidate. Thus the added
weight $g(p,q)$ is determined by the two positions $p$ and $q$.

To maximize covered weight among selections of $t$ candidates ending
at $q$, use the recurrence
\[
 D(t,q)=\max_{p<q}\{D(t-1,p)+g(p,q)\},
 \qquad D(0,0)=0.
\]
Here $D(t-1,p)$ records the largest weight covered by a selection
of $t-1$ candidates whose last position is $p$. The preceding
characterization shows that appending $q$ adds $g(p,q)$, so the recurrence
considers every possible penultimate position and loses no feasible
selection. Taking the maximum over final positions and all $t\le k$
therefore gives the best cover using at most $k$ candidates.
We reserve position $0$ as a \emph{sentinel}, an artificial initial
marker before every interval and outside the selectable candidates.
It supplies the empty selection as the base case; other nonexistent
states have value $-\infty$, so they cannot enter a maximizing
transition. Since there are $O(kM)$ states and each maximization
checks at most $M$ predecessors, this is a polynomial dynamic program.

For circular intervals and $k\ge1$, enumerate a selected candidate
$v$. Credit the demands it hits and delete them. Cutting at $v$
turns the remaining demands into ordinary intervals, on which
the dynamic program selects at most $k-1$ more candidates.
One enumeration contains a candidate of an optimal nonempty
selection, so taking the best is exact. The $k=0$ case is empty.
\end{proof}

\begin{algorithm}[H]
\caption{\IntervalCover: exact minimum cover or budgeted coverage}
\label{alg:interval-cover}
\KwIn{Ordered interval demands with weights; either \emph{minimum} or budget $k$;
ordinary or circular order.}
\KwOut{A minimum hitting set, or a maximum-weight set of size at most $k$.}
\If{the demand family is empty}{\KwRet $\varnothing$\;}
\If{mode is \emph{minimum} and an empty demand exists}{\KwRet infeasible\;}
\If{mode is budget coverage and $k=0$}{\KwRet $\varnothing$\;}
\If{order is circular}{
  Enumerate a forced selected candidate $v$; delete its hit demands and cut there\;
  Solve each residual ordinary instance, with budget $k-1$ in budget mode\;
  \KwRet the best residual solution plus $v$, also considering the empty budget solution\;
}
\eIf{mode is \emph{minimum}}{
  $P\gets\varnothing$\;
  \While{a demand remains}{
    Choose the right endpoint $q$ of an earliest-ending remaining demand\;
    $P\gets P\cup\{q\}$; delete all demands containing $q$\;
  }
  \KwRet $P$\;
}{
  Set $D(0,0)=0$ and all other unreachable states to $-\infty$\;
  Precompute $g(p,q)$: weight of intervals containing $q$ with left endpoint $>p$\;
  \For{$t=1,\ldots,\min\{k,M\}$}{
    \For{$q=1,\ldots,M$}{
      $D(t,q)\gets\max_{p<q}(D(t-1,p)+g(p,q))$; store a maximizing $p$\;
    }
  }
  \KwRet the reconstruction of the best state with $t\le k$, including $(0,0)$\;
}
\end{algorithm}

For minimum circular cover the forced candidate need only range over
one nonempty demand. For budget coverage, empty demands contribute
zero covered weight. The empty family returns the empty set in both
modes. Lemmas~\ref{lem:minimum-interval-cover} and
\ref{lem:budget-coverage} establish exactness.


We now turn these exact budgeted covers into an ordering. Geometrically
increasing budgets let small batches find likely targets early, while
the minimum cover at the final budget ensures that every positive
pattern is eventually found. A shift changes when the sequence crosses
integer budgets; averaging over shifts bounds both positive first-cover
times and the wait on the all-negative route. Since only finitely many
distinct schedules arise, we can then select one deterministically.


\subsection{The batch-ordering guarantee}

Fix a nonempty demand family $U$, its minimum cover size $k_0$, and
the waiting weight $p_0$. The only independent scheduling parameter
is the growth factor $\beta>1$. The functions below express the
shift-averaged bounds on first-cover cost and schedule length. Once
$\beta$ and the current state are fixed, both are determined.

Put
\[
 r_\beta=\lfloor\beta\rfloor,\qquad
 \lambda_\beta=\frac{\beta+1}{2\ln\beta},\qquad
 \nu_\beta=1+\frac1{\ln\beta}.
\]
For integer $t,k\ge1$, define the accounting functions
\[
 \begin{aligned}
 \Phi_\beta(1)&=\frac{r_\beta+1}{2}
                 -\frac{\ln(r_\beta!)}{2\ln\beta},\\
 \Phi_\beta(t)
 &=\Phi_\beta(1)
    +\frac1{2\ln\beta}\int_1^t
       \frac{\lfloor x\rfloor+\lfloor\beta x\rfloor}{x}\,dx,\\
 Z_\beta(k)
 &=k+\frac{k\ln k-\ln(k!)}{\ln\beta}.
 \end{aligned}
\]

We first derive a bound for each shifted schedule. Keeping the integer
budgets in this calculation preserves the savings gained by rounding
down in the early phases, which improves the final approximation factor.

\begin{proposition}[A shifted schedule has a discrete cost bound]
\label{prop:batch-invoice}
For a comparison ordering $Q$, put
$r_Q(x)=\sum_{h:t_Q(h)>x}\pi_h$. The shifted construction $G_u$
using budgets $k^{(j)}=\lfloor\beta^ju\rfloor$ satisfies
\[
\MSSC(G_u)\le\frac{\lfloor u\rfloor+1}{2}w(U)
 +\frac12\sum_{j\ge0}(k^{(j)}+k^{(j+1)})r_Q(\beta^ju).
\]
Each such schedule has length at most
$(1+\beta/(\beta-1))k_0$.
\end{proposition}

A shift that finds likely positive targets early can still leave a long
all-negative route. We therefore average the first-cover cost and the
schedule length under one distribution, so a single shift can balance
both contributions.

\begin{proposition}[One shift distribution bounds both costs]
\label{prop:batch-average}
For the log-uniform shift on $[1,\beta)$,
\[
\mathbb E[\MSSC(G_u)]\le\sum_h\pi_h\Phi_\beta(t_Q(h)),
\qquad \mathbb E[|G_u|]\le Z_\beta(k_0).
\]
They also satisfy $\Phi_\beta(t)\le\lambda_\beta(t-1)+\Phi_\beta(1)$
and $Z_\beta(k)\le\nu_\beta k-1/\ln\beta$.
\end{proposition}

The average guarantees a schedule with a suitably small joint cost.
Because integer budgets change only at explicit breakpoints, a finite
search over the resulting schedules can find one no more costly than
the average.

\begin{proposition}[A finite shift search suffices]
\label{prop:batch-derandomization}
For fixed $\beta$, polynomially many shifts represent every distinct
$G_u$. Minimizing the joint cost
$\MSSC(G_u)+p_0|G_u|$ over these representatives attains at most
its log-uniform average.
\end{proposition}

\begin{proofstrategy}
\begin{enumerate}
\item Exact budget coverage bounds the mass left after each phase.
\item Greedy order within a phase gives a half-batch saving in first-cover cost.
\item Average over the same shift for discovery and waiting, then
      choose the best of finitely many distinct schedules.
\end{enumerate}
The shift averaging establishes the bound; the finite search returns a
deterministic schedule.
\end{proofstrategy}

\begin{lemma}[Geometric-batch ordering]\label{lem:mssc}
For every fixed $\beta>1$, the following guarantee holds.
Let $U$ be a nonempty family of coverable weighted circular intervals.
Let $k_0$ be its minimum cover size and let $p_0\ge0$ be the weight
of a target covered only when the entire schedule finishes.
A deterministic polynomial-time ordering $G$ satisfies, for every
comparison cover ordering $Q$,
\[
 \MSSC(G)+p_0|G|
 \le\sum_{h\in U}\pi_h\Phi_\beta(t_Q(h))+p_0 Z_\beta(k_0).
\]
The simpler bounds
$\Phi_\beta(t)\le\lambda_\beta(t-1)+\Phi_\beta(1)$ and
$Z_\beta(k)\le\nu_\beta k-1/\ln\beta$ also hold.
Every candidate schedule considered by the construction also has
length at most $\bigl(1+\beta/(\beta-1)\bigr)k_0$.
\end{lemma}
\begin{proof}
Proposition~\ref{prop:batch-invoice} bounds each shifted schedule.
Proposition~\ref{prop:batch-average} then bounds their joint average,
and Proposition~\ref{prop:batch-derandomization} selects a schedule
whose joint cost is no larger. Every candidate also satisfies the
length bound, so the selected schedule inherits it.
\end{proof}

\begin{algorithm}[H]
\caption{\BatchOrders: deterministic family of geometric-batch schedules}
\label{alg:batch-orders}
\KwIn{Weighted circular demands and fixed $\beta>1$.}
\KwOut{All schedules needed for deterministic shift minimization.}
Compute a minimum cover $Q_*$ of size $k_0$ using $\IntervalCover$\;
\For{$k=1,\ldots,k_0-1$}{Compute a maximum-weight budget-$k$ set $Q_k$\;}
Find the first $J$ with $\beta^J\ge k_0$\;
Form the breakpoints $1,\beta$ and $k/\beta^j\in[1,\beta)$,
$0\le j\le J$, $1\le k\le k_0$\;
\ForEach{breakpoint below $\beta$ and each midpoint of consecutive breakpoints}{
  $u\gets$ that representative; $Q\gets$ empty order\;
  \For{$j=0,1,\ldots$}{
    $k\gets\lfloor\beta^ju\rfloor$;
    $P\gets Q_k$ if $k<k_0$, and $P\gets Q_*$ otherwise\;
    Greedily order $P$ by newly covered weight, append its productive unused candidates to $Q$\;
    \If{$k\ge k_0$}{stop this shift's phases\;}
  }
  Retain $Q$\;
}
\KwRet the retained family, removing duplicate schedules\;
\end{algorithm}

An empty demand family needs no schedule. Lemma~\ref{lem:mssc}
bounds the family's minimum joint first-cover and waiting cost.
The number of representatives is $O(k_0\log k_0)$ for fixed $\beta$.


\subsection{Proofs of the batch-ordering propositions}

\begin{proofstrategy}
\begin{enumerate}
\item Compare each budgeted batch with the first equally sized prefix
      of an arbitrary cover ordering, and bound the total number of
      scheduled tests by the geometric budget sum.
\item Use greedy order within each batch to bound its first-cover
      contribution using the half-batch saving.
\item Sum these batch costs, then replace batch residuals by the comparison
      ordering's residual mass to obtain the discrete bound.
\end{enumerate}
\end{proofstrategy}

\begin{proof}[Proof of Proposition~\ref{prop:batch-invoice}]

First use Lemma~\ref{lem:minimum-interval-cover} to compute a minimum
cover $Q_*$, and write $k_0=|Q_*|$. For each integer budget $k<k_0$,
Lemma~\ref{lem:budget-coverage} supplies a maximum-weight set
$Q_k$ of at most $k$ candidates. Fixing these sets deterministically
lets each phase use its best coverage set for that budget; the final
phase uses $Q_*$ and therefore covers every demand.

Fix a comparison cover ordering $Q$, and write
\[
 r_Q(x)=\sum_{h:t_Q(h)>x}\pi_h
\]
for the weight whose first-cover index in $Q$ exceeds $x$.
For a shift $u\in[1,\beta)$, use budgets
$k^{(j)}=\lfloor\beta^ju\rfloor$, $j=0,1,\ldots$.
In each phase, process $Q_{k^{(j)}}$ when $k^{(j)}<k_0$ and $Q_*$
otherwise. Order that set greedily by the weight it covers among
the demands still unhit by earlier queries. Omit repeated queries
and those hitting no remaining demand. Stop after the first phase
whose budget reaches $k_0$.
Two global safeguards come first: each phase covers at least as much
weight as a comparison prefix with the same budget, and geometric
growth bounds the total number of tests.

\begin{claim}[Phase coverage and schedule length]
\label{clm:batch-phase-cover}
After a phase with budget $k<k_0$, the weight still unhit is at
most $r_Q(k)$. The complete shifted schedule has length at most
$\bigl(1+\beta/(\beta-1)\bigr)k_0$.
\end{claim}
\begin{proof}
The union of all queries processed through this phase covers every
demand covered by $Q_k$: a query omitted for hitting no remaining
demand has its covered demands already hit by an earlier phase.
By Lemma~\ref{lem:budget-coverage}, the set $Q_k$ covers at least
as much total weight as the first
$k$ queries of $Q$, since those queries are one feasible
at-most-$k$ candidate set. Hence the schedule has hit at least the
weight hit by that comparison prefix, leaving at most $r_Q(k)$.
The final $Q_*$ phase covers all demands.

Let $J$ be the first phase whose budget reaches $k_0$.
If $J=0$, there are no pre-final queries and the full length is at
most $k_0$. Otherwise the pre-final phases use at most
\[
 \sum_{j<J}\lfloor\beta^ju\rfloor
 \le\sum_{j<J}\beta^ju
 <\frac{\beta}{\beta-1}k_0
\]
queries. The first inequality uses $\lfloor x\rfloor\le x$.
For the strict step, $\lfloor\beta^{J-1}u\rfloor<k_0$ and
integrality of $k_0$ imply $\beta^{J-1}u<k_0$. Reversing the
finite geometric sum and bounding it by the infinite sum gives
\[
 \sum_{j<J}\beta^ju
 =\beta^{J-1}u\sum_{r=0}^{J-1}\beta^{-r}
 <k_0\sum_{r\ge0}\beta^{-r}
 =\frac{\beta}{\beta-1}k_0.
\]
The final minimum cover uses $k_0$ queries, proving the length
bound.
\end{proof}

Phase-end coverage alone would charge every new hit at the end of its
batch. Greedy places the largest new gains first, giving the half-batch
saving captured by the next claim.

\begin{claim}[First-cover cost within one batch]
\label{clm:batch-first-cover}
Suppose a phase processes $m$ queries, newly hitting weights
$g_1,\ldots,g_m$ in order, and starts and ends with unhit weights
$r$ and $r'$. Its contribution to the first-cover sum is at most
\[
 mr'+\sum_{i=1}^m i g_i
 \le mr'+\frac{m+1}{2}(r-r')
 =\frac{m+1}{2}r+\frac{m-1}{2}r'.
\]
\end{claim}
\begin{proof}
After each greedy choice, the remaining coverage of every unused
candidate can only decrease. Therefore the newly hit weights are
nonincreasing: $g_1\ge\cdots\ge g_m$. Pairing ranks $i$ and
$m+1-i$ shows
\[
 \sum_{i=1}^m i g_i\le\frac{m+1}{2}\sum_{i=1}^m g_i.
\]
Indeed, in each pair the earlier rank has at least as much weight
as the later one, so replacing both ranks by their mean cannot
decrease the weighted sum. Explicitly, with $a=(m+1)/2$ and $i<a$,
\[
 (i-a)g_i+(m+1-i-a)g_{m+1-i}
 =(i-a)(g_i-g_{m+1-i})\le0.
\]
The middle term, when $m$ is odd, already has rank $a$.
Summing these pair inequalities proves the displayed bound.
A demand first hit at rank $i$ remains
unhit for $i$ queries in this batch, while each demand
still unhit at the end remains unhit for all $m$ queries. Since
$\sum_i g_i=r-r'$, their combined contribution gives the claim.
\end{proof}

Now let $r_j$ be the unhit weight after phase $j$, and put
$r_{-1}=w(U)$. If phase $j$ processes $m_j$ queries, the preceding
claim bounds its first-cover contribution by
\[
 \frac{m_j+1}{2}r_{j-1}+\frac{m_j-1}{2}r_j.
\]
Let $J$ be the final phase. Since $r_J=0$, adding the batch costs
of Claim~\ref{clm:batch-first-cover} and collecting the two
appearances of each intermediate residual gives
\[
\begin{aligned}
 \MSSC(G_u)
 &\byref{Claim}{clm:batch-first-cover}{\le}\sum_{j=0}^{J}
       \left(\frac{m_j+1}{2}r_{j-1}+\frac{m_j-1}{2}r_j\right)\\
 &=\frac{m_0+1}{2}w(U)
     +\frac12\sum_{j=0}^{J-1}(m_j+m_{j+1})r_j.
\end{aligned}
\]
Here the coefficient of $r_j$ is
$(m_j-1)/2+(m_{j+1}+1)/2=(m_j+m_{j+1})/2$.
All residuals are nonnegative, so $m_j\le k^{(j)}$ permits replacing
each batch length by its budget. Claim~\ref{clm:batch-phase-cover} then bounds
$r_j\le r_Q(k^{(j)})$, giving
\[
 \MSSC(G_u)\le
 \frac{\lfloor u\rfloor+1}{2}w(U)
 +\frac12\sum_{j\ge0}
       \bigl(k^{(j)}+k^{(j+1)}\bigr)r_Q(\beta^ju).
\]
For every pre-final phase, Claim~\ref{clm:batch-phase-cover} gives
$r_j\le r_Q(k^{(j)})=r_Q(\beta^ju)$; the equality uses the
integer-valued comparison first-cover indices. The final phase
has residual zero. The displayed infinite sum only adds
nonnegative upper bounds for the final and later phases.

\end{proof}

\begin{proofstrategy}
\begin{enumerate}
\item Average the discrete cost bound under the log-uniform shift; compute
      its initial floor term and combine the scale intervals into one
      integral, which yields $\Phi_\beta$.
\item Apply $\lfloor x\rfloor\le x$ for the simpler linear bound.
\item Average the same shift's earlier batch lengths, evaluate the
      floor integral by a factorial sum, and use the integral lower bound
      on $\ln(k!)$ to obtain the stated bound on $Z_\beta$.
\end{enumerate}
\end{proofstrategy}

\begin{proof}[Proof of Proposition~\ref{prop:batch-average}]
A logarithmic shift moves the phase boundaries across every scale
with the same density. After substitution, their averaging intervals
tile $[1,\infty)$; this yields one integral for discovery and another
for waiting. Keeping floors in both integrals preserves the discrete savings.

\begin{claim}[Shift averaging yields the discrete first-cover function]
\label{clm:batch-average-first-cover}
The same log-uniform shift satisfies
$\mathbb E[\MSSC(G_u)]\le\sum_h\pi_h\Phi_\beta(t_Q(h))$.
\end{claim}
\begin{proof}

Average the cost bound of Proposition~\ref{prop:batch-invoice} over
a \emph{log-uniform shift}, with
probability density $1/(u\ln\beta)$ on $[1,\beta)$. This density
integrates to one because $\int_1^\beta du/u=\ln\beta$.
Splitting the first term's integral at its integers
gives
\[
 \begin{aligned}
 \mathbb E[\lfloor u\rfloor]
 &=\frac{\displaystyle
     \sum_{k=1}^{r_\beta-1}k\ln\frac{k+1}{k}
        +r_\beta\ln\frac{\beta}{r_\beta}}{\ln\beta}\\
 &=r_\beta-\frac{\ln(r_\beta!)}{\ln\beta}.
 \end{aligned}
\]
For the last equality, expand the finite sum:
\[
 \sum_{k=1}^{r_\beta-1}k(\ln(k+1)-\ln k)
 =(r_\beta-1)\ln r_\beta-\ln((r_\beta-1)!).
\]
Adding $r_\beta(\ln\beta-\ln r_\beta)$ leaves
$r_\beta\ln\beta-\ln(r_\beta!)$. This also holds for $r_\beta=1$,
when the sum is empty. Hence the initial term averages to
$\tfrac12(\mathbb E[\lfloor u\rfloor]+1)w(U)=\Phi_\beta(1)w(U)$.
The intervals $[\beta^j,\beta^{j+1})$ partition $[1,\infty)$,
and substituting $x=\beta^ju$ in each phase gives
\[
\begin{aligned}
 &\sum_{j\ge0}\int_1^\beta
 \frac{\lfloor\beta^ju\rfloor+\lfloor\beta^{j+1}u\rfloor}{u}
 r_Q(\beta^ju)\,du\\
 &\qquad=\sum_{j\ge0}\int_{\beta^j}^{\beta^{j+1}}
 \frac{\lfloor x\rfloor+\lfloor\beta x\rfloor}{x}r_Q(x)\,dx\\
 &\qquad=\int_1^\infty
 \frac{\lfloor x\rfloor+\lfloor\beta x\rfloor}{x}r_Q(x)\,dx.
\end{aligned}
\]
Only finitely many terms have positive residual, since
$r_Q(x)=0$ for $x\ge |Q|$. Dividing by the shift density's
normalizing factor $2\ln\beta$ now yields
\[
 \begin{aligned}
 \mathbb E[\MSSC(G_u)]
 &\le(\Phi_\beta(1))w(U)\\
 &\quad+\frac1{2\ln\beta}\int_1^\infty
      \frac{\lfloor x\rfloor+\lfloor\beta x\rfloor}{x}
      r_Q(x)\,dx\\
 &=\sum_{h\in U}\pi_h\Phi_\beta(t_Q(h)).
 \end{aligned}
\]
For each real $x$, write
$r_Q(x)=\sum_{h:t_Q(h)>x}\pi_h$. Interchanging this finite sum
with the integral makes each hypothesis contribute the integral
from $1$ to its integer first-cover time. Adding the initial
$(\Phi_\beta(1))w(U)$ term therefore gives
$\sum_{h\in U}\pi_h\Phi_\beta(t_Q(h))$.
Finally $\lfloor x\rfloor+\lfloor\beta x\rfloor\le(1+\beta)x$,
so the definition of $\Phi_\beta$ gives
\[
 \Phi_\beta(t)-\Phi_\beta(1)
 \le\frac1{2\ln\beta}\int_1^t(1+\beta)\,dx
 =\lambda_\beta(t-1).
\]
The inequality replaces each floor by its argument, and the
equality uses $\lambda_\beta=(1+\beta)/(2\ln\beta)$.
Using this simpler bound discards the savings from later integer budgets.

\end{proof}

The first-cover sum measures the cost for targets that receive a
positive answer. To include the all-negative route as well, we average
the schedule length under the same shift; the two bounds can then be
combined.

\begin{claim}[Shift-averaged schedule length]
\label{clm:batch-average-length}
Under the same shift distribution,
\[
 \mathbb E[|G_u|]\le Z_\beta(k_0)
 \le\nu_\beta k_0-\frac1{\ln\beta}.
\]
\end{claim}
\begin{proof}
The final cover contributes $k_0$ queries. The earlier phases
contribute at most
$\sum_{j:\,\beta^ju<k_0}\lfloor\beta^ju\rfloor$.
Under the density $1/(u\ln\beta)$, their expected length is
\[
\begin{aligned}
 &\frac1{\ln\beta}\sum_{j\ge0}\int_1^\beta
       \mathbf1[\beta^ju<k_0]\frac{\lfloor\beta^ju\rfloor}{u}\,du\\
 &\qquad=\frac1{\ln\beta}\sum_{j\ge0}
       \int_{\beta^j}^{\beta^{j+1}}
       \mathbf1[x<k_0]\frac{\lfloor x\rfloor}{x}\,dx\\
 &\qquad=\frac1{\ln\beta}\int_1^{k_0}\frac{\lfloor x\rfloor}{x}\,dx.
\end{aligned}
\]
As in Claim~\ref{clm:batch-average-first-cover}, the first
equality uses $x=\beta^ju$, so $du/u=dx/x$; the second joins
the disjoint scale intervals and removes the part where the
indicator is zero.
On each interval $[j,j+1)$, $\lfloor x\rfloor=j$, so summing the
integrals gives
\[
\begin{aligned}
 \int_1^{k_0}\frac{\lfloor x\rfloor}{x}\,dx
 &=\sum_{j=1}^{k_0-1}j(\ln(j+1)-\ln j)\\
 &=(k_0-1)\ln k_0-\ln((k_0-1)!)\\
 &=k_0\ln k_0-\ln(k_0!).
\end{aligned}
\]
The middle equality cancels consecutive logarithmic terms;
the last uses $\ln(k_0!)=\ln k_0+\ln((k_0-1)!)$.
Consequently $\mathbb E[|G_u|]\le Z_\beta(k_0)$. Since $\ln x$
is increasing,
\[
 \ln(k_0!)=\sum_{j=2}^{k_0}\ln j
 \ge\int_1^{k_0}\ln x\,dx
 =k_0\ln k_0-k_0+1.
\]
Indeed, on $[j-1,j]$ the increasing function $\ln x$ is at
most $\ln j$, so summing those integral comparisons gives
the inequality. Evaluating the antiderivative $x\ln x-x$
gives the equality.
Consequently
\[
 Z_\beta(k_0)
 \le k_0+\frac{k_0-1}{\ln\beta}
 =\nu_\beta k_0-\frac1{\ln\beta}.
\]
The formula also gives
$Z_\beta(1)=1$.
\end{proof}

Claims~\ref{clm:batch-average-first-cover}
and~\ref{clm:batch-average-length} hold under the same shift
distribution. Therefore linearity of expectation and $p_0\ge0$ give
\[
 \mathbb E[\MSSC(G_u)+p_0|G_u|]
 \le\sum_{h\in U}\pi_h\Phi_\beta(t_Q(h))
      +p_0Z_\beta(k_0).
\]

\end{proof}

\begin{proof}[Proof of Proposition~\ref{prop:batch-derandomization}]
No random shift is used by the algorithm. Let $J$ be the first
integer with $\beta^J\ge k_0$. All budget changes occur at the
breakpoints $u=k/\beta^j$ for
$0\le j\le J$, $1\le k\le k_0$, inside $[1,\beta)$.
Enumerate these breakpoints and one interior point from each
consecutive cell, construct the resulting orderings, and
take the one minimizing the rational quantity
$\MSSC(G_u)+p_0|G_u|$. There are
$O(k_0)$ distinct intervals: for a fixed integer $k$, a breakpoint
$k/\beta^j$ belongs to $[1,\beta)$ exactly when
$\beta^j\le k<\beta^{j+1}$. These half-open ranges partition
$[1,\infty)$, so at most one exponent qualifies for each $k$.
Adding the boundary shifts and one midpoint per gap still leaves
$O(k_0)$ representatives.
Also $J=O_\beta(1+\log k_0)$. Every constructed ordering
uses distinct candidates and hence has length at most $|V|$.
Evaluate the joint score $\MSSC(G_u)+p_0|G_u|$ for each representative
and return the schedule with minimum score. The breakpoint and midpoint
construction leaves only $O(k_0)$ schedules to compare.
The budgets, selected sets, and their greedy internal orders are
constant on each interval because their only dependence on $u$
is through the integer budgets, with fixed tie-breaking thereafter.
Hence the finite family represents every schedule occurring in
the average of Proposition~\ref{prop:batch-average}: within each cell,
the budget vector, selected sets, and deterministic orders agree with
the representative's. Thus every schedule in the average has one of
the listed joint scores, and the minimum listed score is at most their
weighted average. The waiting-time weight $p_0$ is used
only to choose the schedule; all queries remain supplied candidates.
\end{proof}

